\chapter{Macros}
\index{macros}\index{syntax-rules}
\label{ch:macros}

Macros are what make Lisp languages uniquely powerful. They let you extend the language itself---defining new syntactic forms that look and behave exactly like built-in ones. If you find yourself writing the same boilerplate pattern repeatedly, a macro can eliminate it.

Scheme's macro system is \emph{hygienic}\index{hygienic macros}: variables introduced by a macro never accidentally capture variables from the surrounding code. This eliminates an entire class of bugs that plague macro systems in C and Common Lisp.

This chapter covers \texttt{define-syntax} and \texttt{syntax-rules}---the pattern-based macro system specified by R7RS---and, at the end, Kaappi's procedural alternative, \texttt{er-macro-transformer}.

\section{What Is a Macro?}

A macro is a compile-time transformation rule\index{macro expansion}. When Kaappi encounters a macro call, it does not evaluate the arguments---instead, it rewrites the entire expression according to the macro's rules, and then evaluates the result.

One boundary is worth stating up front: macros transform \emph{expressions}, the tree of parenthesized forms the reader has already parsed. They cannot change how characters are read in the first place. Some Lisps let you extend the reader itself with new lexical syntax; Kaappi (like standard R7RS) does not---the token syntax of Chapter~\ref{ch:first-steps} is fixed.\index{reader}

Consider \texttt{when}:

\begin{lstlisting}[style=repl]
> (when (> 5 3)
    (display "yes")
    (newline))
yes
\end{lstlisting}

\texttt{when} is not a built-in---it is a macro that expands to an \texttt{if}:

\begin{lstlisting}
(if (> 5 3)
    (begin
      (display "yes")
      (newline)))
\end{lstlisting}

The macro takes the \emph{source code} and transforms it into different source code before evaluation:

\begin{center}
\begin{tikzpicture}[
    box/.style={rectangle, draw, rounded corners, font=\ttfamily\scriptsize,
                text width=5.8cm, align=left, inner sep=4pt},
    arr/.style={->, >=stealth, thick},
    lbl/.style={font=\sffamily\scriptsize, fill=white, inner sep=1pt}
]
\node[box] (src) {(when (> 5 3)\\~~(display "yes")\\~~(newline))};
\node[box, below=1.2cm of src] (exp) {(if (> 5 3)\\~~(begin\\~~~~(display "yes")\\~~~~(newline)))};
\draw[arr] (src) -- (exp) node[lbl, midway] {macro expansion};
\end{tikzpicture}
\end{center}

\begin{note}[Coming from Python]
Python has no equivalent feature. Decorators modify functions but cannot create new syntax. Macros operate at the level of code itself---they generate code from code.
\end{note}

\section{Your First Macro}

Define a macro with \texttt{define-syntax}\index{define-syntax} and \texttt{syntax-rules}:

\begin{lstlisting}[style=repl]
> (define-syntax my-when
    (syntax-rules ()
      ((my-when test body ...)
       (if test (begin body ...)))))
> (my-when (> 5 3)
    (display "hello")
    (newline))
hello
\end{lstlisting}

Let's break this down:

\begin{itemize}
    \item \texttt{define-syntax} names the macro (\texttt{my-when}).
    \item \texttt{syntax-rules ()} introduces the transformation rules. The \texttt{()} is a list of literal keywords (empty here).
    \item \texttt{(my-when test body ...)} is the \emph{pattern}---it matches the macro call.
    \item \texttt{(if test (begin body ...))} is the \emph{template}---what the macro expands to.
    \item \texttt{...} (ellipsis) means ``zero or more of the preceding element.''
\end{itemize}

\section{Seeing the Expansion}
\index{macro expansion!inspecting}\index{expand REPL command@\texttt{,expand} REPL command}

When a macro misbehaves, look at what it expands to. The Kaappi REPL's \texttt{,expand} command shows the expansion of any expression without evaluating it:

\begin{lstlisting}[style=repl]
> ,expand (when #t (display "yes"))
(if #t (begin (display "yes")))
> ,expand (my-when (> 5 3) (display "hello") (newline))
(if (> 5 3) (begin (display "hello") (newline)))
\end{lstlisting}

Make \texttt{,expand} a habit while working through this chapter: after each macro you write, expand a sample call and check that the generated code is what you intended. It turns macro debugging from guesswork into reading. The same view exists outside the REPL: \texttt{kaappi expand file.scm}\index{kaappi expand} prints a whole program after macro expansion, which is handy when a macro misbehaves only in the middle of a real file.

\section{Pattern Matching}

\texttt{syntax-rules} patterns work by structural matching:

\begin{itemize}
    \item A bare identifier (like \texttt{test} or \texttt{body}) matches any single expression and binds it as a \emph{pattern variable}.
    \item An identifier followed by \texttt{...} matches zero or more expressions.
    \item Literal parentheses match literal parentheses in the input.
    \item Identifiers listed in the literals list are matched by name, not as pattern variables.
    \item The underscore \texttt{\_}\index{underscore pattern} is a wildcard: it matches any single form without binding it. Because the macro name in a call always matches trivially, real-world macros conventionally write the name position as \texttt{\_}: \texttt{((\_ test body ...) ...)}. This book spells the name out for readability, but expect \texttt{\_} in code you read elsewhere.
\end{itemize}

\subsection{Multiple Clauses}

A macro can have multiple patterns. The first matching pattern wins:

\begin{lstlisting}[style=repl]
> (define-syntax my-and
    (syntax-rules ()
      ((my-and) #t)
      ((my-and x) x)
      ((my-and x rest ...)
       (if x (my-and rest ...) #f))))
> (my-and)
#t
> (my-and 1 2 3)
3
> (my-and 1 #f 3)
#f
\end{lstlisting}

This is a recursive macro: the third clause expands to another call to \texttt{my-and} with fewer arguments, until it reaches a base case.

If no clause matches a call, expansion fails with a compile-time syntax error at the call site. The \texttt{dotimes} macro later in this chapter shows how to catch a misuse deliberately with \texttt{syntax-error}.

\subsection{Literal Keywords}
\label{sec:literal-keywords}

The first argument to \texttt{syntax-rules} is a list of identifiers to match literally rather than as pattern variables:

\begin{lstlisting}[style=repl]
> (define-syntax my-if
    (syntax-rules (then else)
      ((my-if test then consequent else alternate)
       (cond (test consequent)
             (else alternate)))))
> (my-if (> 3 2) then "yes" else "no")
"yes"
\end{lstlisting}

Here, \texttt{then} and \texttt{else} are literal keywords---they must appear in the call exactly as written, not as variable names.

\section{The Ellipsis (\texttt{...})}
\index{ellipsis}

The ellipsis is the key to writing macros that accept variable numbers of subforms.

In a pattern, \texttt{expr ...} matches zero or more expressions. In a template, \texttt{expr ...} replicates the template for each matched element:

\begin{lstlisting}[style=repl]
> (define-syntax my-list
    (syntax-rules ()
      ((my-list) '())
      ((my-list elem rest ...)
       (cons elem (my-list rest ...)))))
> (my-list 1 2 3)
(1 2 3)
\end{lstlisting}

A more useful example---a macro that returns the value of the first expression while evaluating the rest for side effects:

\begin{lstlisting}[style=repl]
> (define-syntax begin0
    (syntax-rules ()
      ((begin0 first rest ...)
       (let ((result first))
         rest ...
         result))))
> (begin0 (+ 1 2)
    (display "side effect")
    (newline))
side effect
3
\end{lstlisting}

\subsection{Nested Ellipsis}

Ellipsis can be nested for macros that process lists of lists:

\begin{lstlisting}[style=repl]
> (define-syntax my-let
    (syntax-rules ()
      ((my-let ((var init) ...) body ...)
       ((lambda (var ...) body ...) init ...))))
> (my-let ((x 5) (y 10))
    (+ x y))
15
\end{lstlisting}

This \texttt{my-let} macro transforms into a lambda application---exactly how \texttt{let} is defined in the language.

\section{Hygiene}
\label{sec:hygiene}
\index{hygienic macros}

Scheme macros are hygienic: temporary variables introduced by a macro cannot clash with variables in the calling code. Consider:

\begin{lstlisting}[style=repl]
> (define-syntax swap!
    (syntax-rules ()
      ((swap! a b)
       (let ((tmp a))
         (set! a b)
         (set! b tmp)))))
> (let ((tmp 100) (x 1) (y 2))
    (swap! x y)
    (list tmp x y))
(100 2 1)
\end{lstlisting}

The \texttt{tmp} inside the macro and the \texttt{tmp} in the calling code are different bindings---the macro's \texttt{tmp} is automatically renamed to avoid capture\index{variable capture}. In an unhygienic system (like C's \texttt{\#define}), this would be a bug.

There is no magic here, and \texttt{,expand} lets you watch the mechanism work. During expansion, every identifier that comes from the \emph{template} (rather than from the macro call) is renamed to a fresh, globally unique name:

\begin{lstlisting}[style=repl]
> ,expand (swap! x y)
(__hyg_1_let ((__hyg_2_tmp x)) (set! x y) (set! y __hyg_2_tmp))
\end{lstlisting}

Read the expansion alongside the template. The template's \texttt{tmp} became \texttt{\_\_hyg\_2\_tmp}, so it can never collide with the \texttt{tmp} at the call site. The \texttt{x} and \texttt{y} came from the macro \emph{call}, not the template, so they keep their names and refer to the caller's bindings. Even \texttt{let} was renamed---the expander renames every template identifier uniformly, and the compiler recognizes renamed references to standard forms---which means the macro still works if the caller has shadowed \texttt{let} with a local variable. (The numbers in the fresh names are a session-wide counter, so yours will differ.)

Hygiene, then, is just consistent renaming: template identifiers are renamed away from the caller's namespace, call-site identifiers pass through untouched. \index{gensym}

\begin{note}[Why hygiene matters]
In C, the macro \texttt{\#define SWAP(a,b) \{int tmp=a; a=b; b=tmp;\}} breaks if you write \texttt{SWAP(tmp, x)} because the user's \texttt{tmp} and the macro's \texttt{tmp} collide. Scheme's system makes this impossible.
\end{note}

\section{Practical Macros}

\subsection{A \texttt{unless} Macro}

Both \texttt{when} and \texttt{unless} are built-in, but they are defined as macros internally. Here is how \texttt{unless} could be implemented:

\begin{lstlisting}
(define-syntax unless
  (syntax-rules ()
    ((unless test body ...)
     (if (not test) (begin body ...)))))
\end{lstlisting}

\subsection{A \texttt{while} Loop}
\index{while}

Scheme does not have a built-in \texttt{while}, but you can add one:

\begin{lstlisting}[style=repl]
> (define-syntax while
    (syntax-rules ()
      ((while test body ...)
       (let loop ()
         (when test
           body ...
           (loop))))))
> (let ((i 0))
    (while (< i 5)
      (display i) (display " ")
      (set! i (+ i 1))))
0 1 2 3 4
\end{lstlisting}

\subsection{A \texttt{dotimes} Loop}
\index{dotimes}

\begin{lstlisting}[style=repl]
> (define-syntax dotimes
    (syntax-rules ()
      ((dotimes (var count) body ...)
       (let loop ((var 0))
         (when (< var count)
           body ...
           (loop (+ var 1)))))))
> (dotimes (i 5)
    (display (* i i))
    (display " "))
0 1 4 9 16
\end{lstlisting}

A call that does not match the pattern---say \texttt{(dotimes 5 ...)}, forgetting the parentheses around \texttt{(i 5)}---fails with a syntax error at expansion time. You can make such misuse explicit by adding a catch-all clause that expands to \texttt{syntax-error}\index{syntax-error}, an R7RS form that signals an error during expansion:

\begin{lstlisting}
(define-syntax dotimes
  (syntax-rules ()
    ((dotimes (var count) body ...)
     (let loop ((var 0))
       (when (< var count) body ... (loop (+ var 1)))))
    ((dotimes other ...)
     (syntax-error "dotimes: expected (variable count)"))))
\end{lstlisting}

The catch-all clause matches anything the first clause rejected, and the \texttt{syntax-error} in its template stops compilation. The message string also documents, right in the macro definition, what shape of call was expected.

\subsection{An Assertion Macro}

\begin{lstlisting}[style=repl]
> (define-syntax assert
    (syntax-rules ()
      ((assert expr)
       (unless expr
         (error "Assertion failed" 'expr)))))
> (assert (= (+ 1 1) 2))
> (assert (= (+ 1 1) 3))
error[KP3000]: Assertion failed (= (+ 1 1) 3)
\end{lstlisting}

The failing assertion raises an error object whose message is \texttt{"Assertion failed"} and whose irritant is the \emph{source code} of the assertion---\texttt{'expr} in the template quotes it, which is why the report shows the exact expression that failed. A handler can read both parts back; Chapter~\ref{ch:errors} shows how.

\subsection{A Timing Macro}

This macro reports how long an expression takes to evaluate. \texttt{current-second}, from the standard \texttt{(scheme time)} library (Chapter~\ref{ch:stdlib}), returns the current time in seconds:

\begin{lstlisting}
(define-syntax time-it
  (syntax-rules ()
    ((time-it expr)
     (let ((start (current-second)))
       (let ((result expr))
         (let ((elapsed (- (current-second) start)))
           (display "Elapsed: ")
           (display elapsed)
           (display "s")
           (newline)
           result))))))
\end{lstlisting}

\section{Local Macros with \texttt{let-syntax}}
\index{let-syntax}\index{letrec-syntax}

You can define macros with limited scope using \texttt{let-syntax}:

\begin{lstlisting}[style=repl]
> (let-syntax ((double (syntax-rules ()
                         ((double x) (+ x x)))))
    (double 21))
42
\end{lstlisting}

The macro \texttt{double} exists only within the body. For mutually recursive macros, use \texttt{letrec-syntax}.

\begin{warning}[Macro arguments may be evaluated more than once]
A \texttt{syntax-rules} macro substitutes its arguments as \emph{code}, not as computed values. If a pattern variable appears several times in the template, its argument expression runs once for each occurrence. A \texttt{square} macro whose body is \texttt{(* x x)} expands \texttt{(square (read))} into \texttt{(* (read) (read))}, reading its input twice. When an argument may have side effects or is expensive, bind it once with \texttt{let} inside the macro before reusing it.
\end{warning}

\section{Procedural Macros: Explicit Renaming}
\label{sec:er-macros}
\index{er-macro-transformer}\index{explicit renaming}\index{procedural macros}\index{SRFI-211}

Everything so far has been \texttt{syntax-rules}: a pattern on one side, a template on the other, and the expander does the rest. That covers most macros, and it is the only macro system R7RS-small specifies. Sometimes, though, the expansion has to be \emph{computed}---checked for shape, built up one element at a time, or generated from a table---and no template can express that. For those cases Kaappi ships a second macro system: \texttt{er-macro-transformer}, the \emph{explicit renaming} transformer standardized by SRFI-211 and imported as \texttt{(srfi 211 explicit-renaming)}. It is an extension beyond R7RS-small, and it is \emph{not} \texttt{syntax-case}\index{syntax-case}, the R6RS system, which Kaappi does not implement.

An explicit-renaming macro is an ordinary procedure that runs at expansion time. It receives three arguments: \texttt{form}, the whole macro call as plain list data; \texttt{rename}, a procedure that turns a symbol into a hygienic identifier; and \texttt{compare}, a procedure that tests whether two identifiers mean the same thing. Whatever the procedure returns \emph{is} the expansion, so quasiquote (Chapter~\ref{ch:lists}) is the natural tool for building it. Here is \texttt{swap!} again, written by hand:

\begin{lstlisting}[style=repl]
> (import (srfi 211 explicit-renaming))
> (define-syntax swap!
    (er-macro-transformer
      (lambda (form rename compare)
        (let ((a (cadr form))
              (b (caddr form))
              (tmp (rename 'tmp)))
          `(let ((,tmp ,a))
             (set! ,a ,b)
             (set! ,b ,tmp))))))
> (let ((tmp 100) (x 1) (y 2))
    (swap! x y)
    (list tmp x y))
(100 2 1)
\end{lstlisting}

Set this beside the \texttt{syntax-rules} version in Section~\ref{sec:hygiene}. There the template's \texttt{tmp} was renamed automatically; here \texttt{(rename 'tmp)} does the same job explicitly, producing the kind of fresh identifier the \texttt{,expand} output showed. The pieces that came from the call, \texttt{a} and \texttt{b}, are spliced in untouched, exactly as pattern variables are. Renaming the same symbol twice within one expansion yields the same identifier, so a binding site and its references agree. A fully defensive version would also write \texttt{(rename 'let)} and \texttt{(rename 'set!)}, so that the macro survives a caller who has shadowed those names; \texttt{syntax-rules} did that for every template identifier without being asked, and that is the convenience you give up here.

\begin{note}[The transformer runs at expansion time]
The \texttt{lambda} inside \texttt{er-macro-transformer} is evaluated when the \texttt{define-syntax} is, and called once per use of the macro---before any of the surrounding code runs. It cannot see runtime variables, and an \texttt{error} raised inside it is reported as a syntax error at the macro's use site rather than as a runtime exception.
\end{note}

\subsection{\texttt{compare}: Matching Keywords}

\texttt{compare} answers the question a literal in Section~\ref{sec:literal-keywords} answers: do these two identifiers refer to the same binding? Hand it a piece of the input and a renamed keyword:

\begin{lstlisting}[style=repl]
> (define-syntax my-if
    (er-macro-transformer
      (lambda (form rename compare)
        (let ((test (list-ref form 1))
              (kw1 (list-ref form 2))
              (yes (list-ref form 3))
              (kw2 (list-ref form 4))
              (no (list-ref form 5)))
          (unless (and (compare kw1 (rename 'then))
                       (compare kw2 (rename 'else)))
            (error "my-if: expected then/else keywords"))
          `(if ,test ,yes ,no)))))
> (my-if (> 3 2) then "yes" else "no")
"yes"
\end{lstlisting}

The check is binding-aware, exactly like a \texttt{syntax-rules} literal. If the caller has a local variable named \texttt{else} in scope at the use site, \texttt{(compare kw2 (rename 'else))} returns \texttt{\#f}: the caller's \texttt{else} and the macro's \texttt{else} are different things, and the macro refuses the form just as a \texttt{syntax-rules} version with \texttt{(then else)} as its literals would.

\subsection{Validation with Good Errors}

Because the transformer sees the whole form, it can check the shape and refuse bad input at expansion time, with a message that says what was wrong:

\begin{lstlisting}[style=repl]
> (define-syntax strict-when
    (er-macro-transformer
      (lambda (form rename compare)
        (if (= (length form) 3)
            `(if ,(cadr form) ,(caddr form))
            (error "strict-when: one body form, got"
                   (- (length form) 2))))))
> (strict-when (> 3 2) 'ok)
ok
\end{lstlisting}

Save the import and the definition as \texttt{strict.scm}, add a good use whose result you \texttt{display}, and then a bad one: \texttt{(strict-when \#t 1 2 3)}. Running the file prints the good result and then stops at the bad use, naming the file, the line and column of that use, and your own message with its irritant:

\begin{Verbatim}[fontsize=\footnotesize]
$ kaappi strict.scm
ok
strict.scm:13:1: syntax-error[KP2002]: strict-when: one body form, got 3
\end{Verbatim}

A \texttt{syntax-rules} macro handed a form none of its clauses match can only report \texttt{compile error[KP2001]: invalid syntax}. Here the message is whatever you wrote, and it appears once, at the offending use, before anything runs.

\subsection{When to Reach for It}

Reach for \texttt{syntax-rules} first: it is shorter, declarative, and portable to every R7RS Scheme. Reach for \texttt{er-macro-transformer} when the expansion must be computed---validated with a precise message, built by deciding something per element (the kind of recursion a template cannot express), or generated from data your program holds. The cost is symmetrical: you construct the output \texttt{syntax-rules} would have constructed for you, and you take over renaming every identifier the expansion introduces. Portable code can test for the library itself:

\begin{lstlisting}[style=repl]
> (cond-expand
    ((library (srfi 211 explicit-renaming)) 'procedural)
    (else 'none))
procedural
\end{lstlisting}

\section{When to Use Macros}

Macros are powerful but should not be overused. Use a macro when:

\begin{itemize}
    \item You need to control \emph{when} or \emph{whether} arguments are evaluated (like \texttt{when}, \texttt{unless}, \texttt{and}).
    \item You need to introduce new bindings (like \texttt{let}, \texttt{dotimes}).
    \item You want to generate code at compile time to avoid runtime overhead.
    \item You are building a domain-specific language (DSL) with custom syntax.
\end{itemize}

Do \emph{not} use a macro when a function would suffice. Functions are simpler, composable (you can pass them to \texttt{map}), and easier to debug. If your ``macro'' does not need to delay evaluation or introduce bindings, make it a function.

\section*{Summary}

\begin{itemize}
    \item Macros transform code at expansion time, before evaluation, letting you add new syntactic forms.
    \item \texttt{syntax-rules} matches input against patterns and rewrites it using output templates.
    \item The ellipsis (\texttt{...}) matches and expands sequences of forms and can be nested.
    \item Scheme macros are hygienic, so introduced identifiers never accidentally capture the user's variables.
    \item The REPL's \texttt{,expand} command shows what a call expands to---the first tool to reach for when a macro misbehaves.
    \item \texttt{er-macro-transformer} (SRFI-211) writes the expansion as a procedure: \texttt{rename} supplies hygiene and \texttt{compare} matches keywords. Use it when the expansion must be computed or validated; otherwise stay with \texttt{syntax-rules}.
    \item \texttt{let-syntax} defines local macros; reach for macros only when a function cannot do the job.
\end{itemize}

\section*{Exercises}

\begin{exercise}[Swap!]
Write a \code{swap!} macro that swaps the values of two variables using \code{set!}. Verify that hygiene prevents the macro's temporary variable from clashing with user code.

\begin{lstlisting}[style=repl]
> (let ((a 1) (b 2))
    (swap! a b)
    (list a b))
(2 1)
> (let ((tmp 99) (x 10) (y 20))
    (swap! x y)
    (list tmp x y))
(99 20 10)
\end{lstlisting}
\end{exercise}

\begin{exercise}[Assert-Equal]
Write an \code{assert-equal} macro that takes two expressions, evaluates both, and raises an error showing the source code of each expression and their values if they are not \code{equal?}. (Hint: use \code{'expr} in the template to capture source code.)

\begin{lstlisting}[style=repl]
> (assert-equal (+ 1 2) 3)
;; no output --- assertion passed
> (assert-equal (* 2 3) 7)
error[KP3000]: assertion failed: (* 2 3) evaluated to 6, expected 7
\end{lstlisting}
\end{exercise}

\begin{exercise}[Do-Times with Named Counter]
Write a \code{do-n-times} macro similar to the \code{dotimes} macro from this chapter, but the user supplies the counter variable name, the count, \emph{and} an optional step size (defaulting to 1). You will need two \code{syntax-rules} clauses.

\begin{lstlisting}[style=repl]
> (do-n-times (i 5)
    (display i) (display " "))
0 1 2 3 4
> (do-n-times (i 10 3)
    (display i) (display " "))
0 3 6 9
\end{lstlisting}
\end{exercise}

\begin{exercise}[With-Logging]
Write a \code{with-logging} macro that prints an entry message before evaluating the body and an exit message after, then returns the body's result. The macro should accept a label string as its first argument.

\begin{lstlisting}[style=repl]
> (with-logging "compute"
    (+ 1 2))
[enter: compute]
[exit: compute]
3
\end{lstlisting}

As a bonus, make it print the elapsed time in the exit message using \code{current-second}.
\end{exercise}

\begin{exercise}[Counting at Expansion Time]
Write \code{count-args}, an explicit-renaming macro that expands to the number of subforms it was given. The transformer computes the count, so the expansion is a literal integer.

\begin{lstlisting}[style=repl]
> (count-args a b c)
3
> (count-args)
0
\end{lstlisting}

Use \code{,expand} to confirm that \code{(count-args a b c)} expands to plain \code{3}, then explain why the arguments are never evaluated.
\end{exercise}

In the next chapter (Chapter~\ref{ch:libraries}), we cover the library system---how to organize code into modules with clean interfaces.

\chapterend
